% #############################################################################
% RESUMO em Português
% !TEX root = ../main.tex
% #############################################################################
% reset acronyms
\acresetall

\begingroup
\singlespacing
\setlength{\parskip}{8pt}
\setlength{\parindent}{1.5em}

% use \noindent in firts paragraph
% \noindent

\noindent A regulamentação da União Europeia impõe obrigações organizacionais cujas implicações para cada parte interessada são frequentemente difíceis de identificar. Traduzir estas obrigações em representações das atividades diárias, responsabilidades e dependências exige ligar sistematicamente os requisitos legais à evidência operacional. Esta dissertação propõe um método supervisionado por humanos, apoiado por inteligência artificial generativa, para identificar, especificar e materializar \textit{viewpoints} de arquitetura específicos de cada parte interessada.

Seguindo a DSRM, com base na norma ISO/IEC/IEEE 42010, o método combina entrevistas às partes interessadas, extração do perfil operacional, análise da relevância regulamentar, especificação de \textit{viewpoints} e geração de modelos ArchiMate. A rastreabilidade associa elementos e relações arquiteturais ao texto jurídico vinculativo e à evidência operacional documentada, enquanto a distinção explícita da atribuição de responsabilidades, as etapas de aprovação humana e a reconciliação entre artefactos limitam o conteúdo gerado.

O método é demonstrado com o DORA, através das perspetivas das partes interessadas afetadas. As descrições arquiteturais resultantes relacionam os requisitos regulamentares com o trabalho específico de cada função, distinguindo responsabilidades diretas, dependências externas e responsabilidades de atribuição incerta. As apreciações preliminares das partes interessadas indicam que as \textit{views} formais podem apoiar a compreensão da regulamentação e revelar lacunas no contexto operacional, permitindo aperfeiçoar os \textit{viewpoints}. Identificam também limitações nas narrativas explicativas e nos diagramas simplificados. As demonstrações sustentam a viabilidade da abordagem, com oportunidades de desenvolvimento. O contributo consiste num processo rastreável para comunicar o impacto regulamentar por parte interessada, sem avaliar a conformidade da organização.

\endgroup